\section{Service system and wave-dependent capacities}
\label{sec:capacity}

We first translate local service times into capacities for dated tasks.
These capacities are the inputs to the finite-fleet scheduling model in
Section~\ref{sec:model}. Figure~\ref{fig:service-architecture} distinguishes
the shipment stages from the spatial pooling within a local wave.

\subsection{Service architecture and operational decisions}

Cities $k\in\{1,2\}$ are connected by a bidirectional corridor. Let
$\mathcal D=\{+,-\}$ be the direction set, with
$o(+)=r(-)=1$ and $o(-)=r(+)=2$. The order set is
$\mathcal L=\bigcup_{d\in\mathcal D}\mathcal L^d$, where each class has
one direction. An order class
$\ell\in\mathcal L^d$ travels from $o(d)$ to $r(d)$ and has quantity
$D_\ell\geq0$, release time $s_\ell$, completion deadline $e_\ell$, and
rejection penalty $p_\ell\geq0$ per unit. Quantities are continuous and
divisible: one unit denotes one pickup-stop and one delivery-stop equivalent.
Thus, the formulation does not assert that a particular indivisible parcel
can be split. We use the time nodes $\mathcal T=\{0,\ldots,T\}$,
where $0\leq s_\ell\leq e_\ell\leq T$ are integers and a period lasts
$\Delta t>0$ minutes. Corridor travel takes $\tau\in\mathbb Z_{>0}$
periods, or $\tau_{\mathrm{phys}}=\tau\Delta t$ minutes.

Initially city $k$ contains $N_k$ human-driven vehicles and $A_k$ automated
vehicles, where $N_k,A_k\in\mathbb Z_{\geq0}$. Write $N=N_1+N_2$ and
$A=A_1+A_2$. The respective capacities are
$M>0$ and $M_A>0$. Human-driven vehicles perform local tours and may also
cross the corridor; automated vehicles operate between hubs. A local tour
starts and ends at its city's hub. A three-leg shipment is collected to
the origin hub, transported between hubs, and distributed from the
destination hub. Direct service keeps the shipment on a human-driven
vehicle across cities. The city-inventory model credits that vehicle to
the other city's available fleet at its modeled trip-completion time. Independent empty
human-driven transfers are also permitted, whether or not there is an order
requiring service.

\begin{table}[tbp]
\centering\small
\caption{Main notation and units for finite-horizon scheduling.}
\label{tab:notation}
\begin{tabular}{@{}p{0.25\linewidth}p{0.69\linewidth}@{}}
\toprule
Symbol & Meaning and unit\\
\midrule
\multicolumn{2}{@{}l}{\textit{Sets and inputs}}\\
$\mathcal L^d,\mathcal T$ & Order classes in direction $d$; discrete time nodes\\
$D_\ell,s_\ell,e_\ell$ & Stop-equivalent demand; release and deadline in periods\\
$N_k,A_k;M,M_A$ & Initial human-driven and automated fleets; load limits in stop equivalents\\
$\Delta t,\tau,\tau_E$ & Minutes per period; loaded and empty intercity travel in periods\\
$a_k,b_k,\rho_k$ & In-zone time coefficients; one-way local stem time in minutes\\
$c_H,c_A,c_D,c_E;p_\ell$ & Currency per vehicle-minute; currency per unserved stop equivalent\\
\midrule
\multicolumn{2}{@{}l}{\textit{Wave and scheduling quantities}}\\
$q,v,n=q/v$ & Total wave stops, vehicle count, and average load per vehicle\\
$x,y,w,e$ & Local, automated, direct, and empty vehicle counts on service arcs\\
$u,\widetilde u$ & Binary selectors for local and direct vehicle counts\\
$g,\widehat g,h,z$ & Local, automated, direct, and unserved stop-equivalent quantities\\
\bottomrule
\end{tabular}
\end{table}

\begin{figure}[tbp]
\centering
\setlength{\unitlength}{1mm}
\begin{picture}(146,65)
\put(0,56){\makebox(146,5)[l]{\small\textbf{Service tasks for one order direction}}}
\put(0,43){\framebox(40,9){\small Collection to hub 1}}
\put(53,43){\framebox(40,9){\small AV hub 1 to hub 2}}
\put(106,43){\framebox(40,9){\small Distribution at city 2}}
\put(40,47.5){\vector(1,0){13}}
\put(93,47.5){\vector(1,0){13}}
\put(0,34){\makebox(146,5)[l]{\small Direct: one human-driven vehicle collects, crosses cities, and delivers.}}
\put(0,26){\makebox(146,5)[l]{\small Direct and empty tasks move vehicles between city inventories at completion.}}
\put(0,17){\makebox(146,5)[l]{\small\textbf{Local wave geometry}}}
\put(0,2){\framebox(57,10){\small Total area $S_k$, total stops $q$}}
\put(57,7){\vector(1,0){20}}
\put(77,2){\framebox(69,10){\small $v$ balanced zones: area $S_k/v$, stops $q/v$}}
\end{picture}
\caption{Service architecture and the balanced local-wave approximation.
The direct description specifies tasks and inventory accounting; it does
not specify unresolved hub-to-service-region connections. Reverse-direction
tasks use the same architecture.}
\label{fig:service-architecture}
\end{figure}

\subsection{Local routing approximation and wave geometry}

For $q\geq0$, define
\begin{equation}
 f_k(q)=a_kq+b_k\sqrt q,\qquad a_k>0,\quad b_k\geq0,
 \qquad
 F_k(\theta)=\left(\frac{\sqrt{b_k^2+4a_k\theta}-b_k}{2a_k}\right)^2,
 \quad \theta\geq0.
 \label{mod:local-function}
\end{equation}
The first term represents handling and the second approximates routing time
under the spatial scaling associated with the BHH limit
\cite{Beardwood1959}. The function $F_k$ is the inverse of $f_k$.
Let $\rho_k\geq0$ denote the average one-way stem time between a hub
and a local service region. Our local wave approximation
assumes that a wave of $v$ vehicles shares a total load $q$ across equal-area
zones, giving per-vehicle duration $2\rho_k+f_k(q)/v$.
This zoning rule is an engineering approximation for finite customer
samples, motivated by the asymptotic routing relation rather than a
finite-sample theorem. For $v>0$, the average vehicle load is $n=q/v$;
the wave quantity $q$ and the vehicle count $v$ are distinct.
Quantities $q,D_\ell,M,M_A$ share the same stop-equivalent unit.
The coefficients $a_k$ and $b_k$ have units of minutes per stop equivalent
and minutes per square root of a stop equivalent, respectively;
$\rho_k$ is in minutes, and $f_k(q)$ is aggregate in-zone service time
in minutes. Period durations are converted to physical time by $\Delta t$.

\paragraph{Why the wave quantity enters the routing term.}
For the balanced zoning approximation, let $S_k$ be the fixed service
area and $\bar\beta_k$ a routing-time coefficient incorporating speed
and the distance metric. Assigning each of $v$ vehicles an area $S_k/v$
and $q/v$ uniformly distributed stops gives the per-vehicle in-zone time
\begin{equation}
 a_k\frac qv+\bar\beta_k\sqrt{\frac{S_k}{v}\frac qv}
 =\frac{a_kq+b_k\sqrt q}{v},\qquad
 b_k=\bar\beta_k\sqrt{S_k}.
 \label{eq:wave-geometry}
\end{equation}
This explains the adopted $f_k(q)/v$ expression. By contrast, $f_k(q/v)$
assigns each vehicle's load to the full service area, producing a routing
term $b_k\sqrt{q/v}$. The distinction concerns spatial pooling, not just
load allocation. The area representation explains an aggregate coefficient;
$b_k$ is an input to the model, not an estimate of area or speed reported
here. The separation of stem, in-zone routing, and stop service follows
the continuous-approximation approach illustrated by
\citet{Banerjee2022}, whose fixed single-vehicle zones serve a
different design problem.

\subsection{Service-arc capacities and pooling properties}
\label{subsec:capacity-analysis}

A local wave is indexed by its start, end, city, and order direction,
$(i,j,k,d)$. Waves that share a start but have different end times are
separate waves; their loads are not pooled to obtain an extra routing
economy. Define the mother arc sets
\begin{align}
 L_k&=\left\lceil\frac{2\rho_k+f_k(M)}{\Delta t}\right\rceil,
 &\mathcal A_k&=\{(i,j)\in\mathcal T^2:1\leq j-i\leq L_k\},
 \label{mod:local-arcs}\\
 \mathcal A_A&=\{(i,j)\in\mathcal T^2:j-i=\tau\},
 &\mathcal A_E&=\{(i,j)\in\mathcal T^2:j-i=\tau_E\},
 \label{mod:movement-arcs}\\
 L_D&=\tau+\left\lceil
 \frac{2\max_k\rho_k+f_1(M)+f_2(M)}{\Delta t}\right\rceil,
 &\mathcal A_D&=\{(i,j)\in\mathcal T^2:\tau\leq j-i\leq L_D\}.
 \label{mod:direct-arcs}
\end{align}
Here $\tau_E\in\mathbb Z_{>0}$ is the empty-transfer duration; its default
is $\tau$. These finite domains are part of the dispatch model. In
particular, all tasks finish within the actual operating horizon $T$.

For $v=1,\ldots,N$, precompute the local capacity
\begin{equation}
 K^k_{ij}(v)=\min\{vM,F_k(v[(j-i)\Delta t-2\rho_k]^+)\},
 \qquad [x]^+=\max\{x,0\}.
 \label{mod:local-capacity}
\end{equation}
Fix a local arc and suppress its city index. Write $D=(j-i)\Delta t$,
$B=(D-2\rho)^+$, and $K(v)=\min\{vM,f^{-1}(vB)\}$, extended to
$v>0$, with $K(0)=0$. The load constraint $q/v\le M$ and the duration
constraint $2\rho+f(q)/v\le D$ define the same capacity for positive $q$.

\begin{proposition}[Capacity economies within a fixed service arc]
\label{prop:analysis-capacity}
For fixed $D,\rho,M$, the ratio $K(v)/v$ is nondecreasing for $v>0$.
Consequently,
\begin{equation}
 K(v_1+v_2)\ge K(v_1)+K(v_2),\qquad v_1,v_2\ge0.
 \label{eq:analysis-superadditive}
\end{equation}
These statements apply in particular to integer fleet sizes. They do not
imply that $K$ is convex or that either inequality is strict.
\end{proposition}

Pooling within the same city, direction, start, and end time weakly
increases modeled capacity relative to separate waves of equal duration.
The gain can disappear when load limits bind. The derivation and proof
are in Appendix~\ref{app:proofs}.

\subsection{Direct-service capacity and route conventions}

For a direct wave in direction $d$, the origin collection is zoned, whereas
each vehicle's destination load is spread over the destination service
area. We adopt the following aggregate direct-service duration approximation:
\begin{equation}
 \Psi_d(q,v)=2\rho_{o(d)}+
 \frac{f_{o(d)}(q)}{v}+\tau\Delta t+f_{r(d)}(q/v).
 \label{mod:direct-duration}
\end{equation}
Set $\widetilde K^d_{ij}(v)$ equal to the largest $q\in[0,vM]$ satisfying
$\Psi_d(q,v)\leq(j-i)\Delta t$, or zero if this set is empty. We retain
the convention $\widetilde K^d_{ij}(v)=0$ when $j-i=\tau$.
The table also has a closed form. For
\begin{equation}
 B_{ij}^d=(j-i-\tau)\Delta t-2\rho_{o(d)},\qquad
 \alpha_{dv}=\frac{a_{o(d)}+a_{r(d)}}{v},\qquad
 \beta_{dv}=\frac{b_{o(d)}}{v}+\frac{b_{r(d)}}{\sqrt v},
 \label{mod:direct-coefficients}
\end{equation}
if $j-i>\tau$ and $B_{ij}^d>0$, then
\begin{equation}
 \widetilde K^d_{ij}(v)=\min\left\{vM,
 \left(\frac{\sqrt{\beta_{dv}^2+4\alpha_{dv}B_{ij}^d}
                   -\beta_{dv}}{2\alpha_{dv}}\right)^2\right\};
 \label{mod:direct-capacity}
\end{equation}
otherwise the capacity is zero. Unlike a symmetric pooled-tour formula,
\eqref{mod:direct-duration} generally gives different capacities in the two
directions.


